\section{Preliminaries}
\label{sec:prelim}

\subsection{Inscribed polytopes}

\begin{cited}{Notation}{{\cite[Section~1]{horvath-langi-2016}}}
  Throughout the paper, by a polytope we mean a convex polytope. The vertex set
  of a polytope $P$ is denoted by $V(P)$. We denote the family of
  $d$-dimensional polytopes, with $n$ vertices and inscribed in the unit sphere
  $\mathbb{S}^{d-1}$, by $\mathcal{P}_{d}(n)$. We denote $d$-dimensional
  volume by $\vol_{d}$, and set
  $v_{d}(n) = \max\{\vol_{d}(P) : P \in \mathcal{P}_{d}(n)\}$.
\end{cited}

\begin{remark}[{\cite[Remark~2.1]{medgyes-maxvol-d4}}]
  \label{rem:attained}
  For $n \geq d + 1$, $v_{d}(n)$ is attained by some
  $P \in \mathcal{P}_{d}(n)$.
\end{remark}

\begin{definition}[{\cite[Definition~2.1, Section~2.2 and Proposition~2.16]{ziegler-1995}}]
  A \emph{facet} of a $d$-polytope is a face of dimension $d-1$, and a
  \emph{ridge} is a face of dimension $d-2$. A polytope is \emph{simplicial} if
  every facet is a simplex. Two polytopes are \emph{combinatorially
  equivalent} if their face lattices are isomorphic; an equivalence class is a
  \emph{combinatorial type}.
\end{definition}

Equivalently, two polytopes are combinatorially equivalent if their
facet--vertex incidence matrices differ only by permutations of rows and
columns \cite[Section~15.1.1]{henk-richter-gebert-ziegler-2017}. For
simplicial polytopes, which is the form in which we use the definition, this
means that a bijection between the vertex sets carries the vertex sets of the
facets of one onto those of the other.

\begin{definition}[{\cite[Definition~1]{horvath-langi-2016}}]
  \label{def:propertyz}
  Let $P \in \mathcal{P}_{d}(n)$ be a $d$-polytope with
  $V(P) = \{p_{1}, p_{2}, \ldots, p_{n}\}$. If for each $i$, there is an open
  set $U_{i} \subset \mathbb{S}^{d-1}$ such that $p_{i} \in U_{i}$, and for any
  $q \in U_{i}$, we have
  \begin{equation*}
    \vol_{d}\bigl(\conv((V(P) \setminus \{p_{i}\}) \cup \{q\})\bigr)
      \leq \vol_{d}(P),
  \end{equation*}
  then we say that $P$ satisfies Property Z.
\end{definition}

\begin{remark}[{\cite{berman-hanes-1970}, \cite[Remark~2.4]{medgyes-maxvol-d4}}]
  \label{rem:propertyz}
  Every $P \in \mathcal{P}_{d}(n)$ with $\vol_{d}(P) = v_{d}(n)$ satisfies
  Property Z.
\end{remark}

\begin{lemma}[{\cite[Lemma~1]{horvath-langi-2016}}]
  \label{lem:simplicial}
  Every $P \in \mathcal{P}_{d}(n)$ satisfying Property Z is simplicial.
\end{lemma}

\subsection{Signed volume and certificates}

For $v_{1}, \ldots, v_{d} \in \mathbb{R}^{d}$, $\det(v_{1}, \ldots, v_{d})$
denotes the determinant of the matrix with columns $v_{1}, \ldots, v_{d}$.

\begin{definition}[{\cite[Definition~3.1]{medgyes-maxvol-d4}}]
  \label{def:signed-volume}
  An \emph{orientation} of a finite set is an ordering of its elements up to
  even permutations. Let $\mathcal{C} = (n, \mathcal{F})$, where
  $\mathcal{F}$ is a finite set of ordered $d$-tuples
  $F = (F_{1}, \ldots, F_{d})$ of distinct elements of $\{1, \ldots, n\}$,
  called facets. Deleting $F_{a}$ gives a ridge, with the orientation of
  $(F_{1}, \ldots, \widehat{F_{a}}, \ldots, F_{d})$ if $a$ is even and the
  opposite one if $a$ is odd. The orientation of $\mathcal{C}$ is
  \emph{coherent} if every ridge lies in exactly two facets and these induce
  opposite orientations on it. A point
  $p = (p_{1}, \ldots, p_{n}) \in (\mathbb{R}^{d})^{n}$ \emph{realises}
  $\mathcal{C}$ if $P = \conv\{p_{1}, \ldots, p_{n}\}$ has the $p_{i}$ as its
  vertices and its facets are exactly the simplices
  $\conv\{p_{j} : j \in F\}$, $F \in \mathcal{F}$. The \emph{signed volume}
  of $p$ is
  \begin{equation}
    \label{eq:signed-volume}
    D(p) = \sum_{F \in \mathcal{F}} \det(p_{F_{1}}, \ldots, p_{F_{d}}).
  \end{equation}
\end{definition}

\begin{lemma}[{\cite[Lemma~3.2]{medgyes-maxvol-d4}}]
  \label{lem:volume}
  If the orientation of $\mathcal{C}$ is coherent and $p$ realises
  $\mathcal{C}$, then
  \[
    d!\,\vol_{d}(\conv\{p_{1}, \ldots, p_{n}\}) = |D(p)|.
  \]
\end{lemma}

\begin{definition}[{\cite[Definition~3.4]{medgyes-maxvol-d4}}]
  \label{def:certificate}
  Let $x_{i} = (x_{i,1}, \ldots, x_{i,d})$ be the coordinates of $p_{i}$, and
  $v(x)$ the vector of all monomials of degree at most $2$ in the $x_{i,a}$. A
  \emph{certificate} for $\mathcal{C}$ and $\beta \in \mathbb{Q}$ is a
  symmetric positive definite matrix $Q$ with rational entries together with
  polynomials $h_{1}, \ldots, h_{n}$ with rational coefficients such that
  \begin{equation}
    \label{eq:certificate}
    \beta - D(x) = v(x)^{\mathsf{T}} Q\, v(x)
      + \sum_{i=1}^{n} h_{i}(x)\bigl(|x_{i}|^{2} - 1\bigr)
  \end{equation}
  holds identically as polynomials in the $x_{i,a}$.
\end{definition}

\begin{proposition}[{\cite[Proposition~3.5]{medgyes-maxvol-d4}}]
  \label{prop:certificate}
  Suppose the orientation of $\mathcal{C}$ is coherent and a certificate for
  $\mathcal{C}$ and $\beta$ exists. Then every polytope inscribed in
  $\mathbb{S}^{d-1}$ that realises $\mathcal{C}$ has volume at most
  $\beta/d!$.
\end{proposition}

A certificate is finite rational data, and its hypotheses (the identity
\eqref{eq:certificate}, positive definiteness of $Q$, coherence, and the facet
sets of the type) are checked in exact arithmetic as described in
\cite[Section~3.3]{medgyes-maxvol-d4}.

\subsection{Simplicial \texorpdfstring{$4$}{4}-polytopes with nine vertices}

\begin{theorem}[{\cite{altshuler-bokowski-steinberg-1980},
    \cite[Table~3]{firsching-2017}}]
  \label{thm:count}
  There are exactly $1142$ combinatorial types of simplicial $4$-polytopes with
  $9$ vertices.
\end{theorem}

The count is due to Altshuler, Bokowski and Steinberg; table numbers of
\cite{firsching-2017} refer to its arXiv version.
We use the realisations of all $1142$ types given there (Lemma~\ref{lem:data}).

\subsection{Direct sums}

\begin{cited}{Direct sum}{{\cite[Section~15.1.3]{henk-richter-gebert-ziegler-2017}}}
  Let $P \subseteq \mathbb{R}^{d}$ be a $d$-polytope with $n$ vertices and $m$
  facets, and $P' \subseteq \mathbb{R}^{d'}$ a $d'$-polytope with $n'$
  vertices and $m'$ facets, both with the origin in their interiors. The
  \emph{direct sum}
  $P \oplus P' = \conv(\{(p, 0) : p \in P\} \cup \{(0, p') : p' \in P'\})$ is
  $(d + d')$-dimensional, and has $n + n'$ vertices and $m \cdot m'$ facets.
\end{cited}

We use this for two regular polygons in orthogonal $2$-planes through the
origin, in Lemma~\ref{lem:freesum}. A polygon has as many edges (its facets)
as vertices, so for polygons with $m$ and $l$ vertices the direct sum has
$m + l$ vertices and $ml$ facets. For $m \geq 3$, $R_{m}$ denotes a regular
$m$-gon inscribed in a unit circle; $\area(R_{m}) = \frac{m}{2}\sin\frac{2\pi}{m}$.

\subsection{Bivectors}

Let $e_{1}, \ldots, e_{N}$ be the standard basis of $\mathbb{R}^{N}$ and
$v \cdot w$ the standard inner product. We use the usual identification of
bivectors with skew-symmetric matrices: the space $\Lambda^{2}\mathbb{R}^{N}$
of \emph{bivectors} is the space of real skew-symmetric $N \times N$ matrices,
and for $v, w \in \mathbb{R}^{N}$ the \emph{wedge product} is
\[
  v \wedge w = vw^{\mathsf{T}} - wv^{\mathsf{T}}.
\]
It is bilinear and satisfies $w \wedge v = -\,v \wedge w$, so $v \wedge v = 0$.
A skew-symmetric matrix $X$ is determined by its entries $X_{kl}$, $k < l$,
above the diagonal, and $e_{kl} = e_{k} \wedge e_{l}$ has entry $1$ at
$(k, l)$, $-1$ at $(l, k)$ and $0$ elsewhere. Hence the $e_{kl}$, $k < l$, form
a basis of $\Lambda^{2}\mathbb{R}^{N}$, we write
$X = \sum_{k<l} X_{kl}\, e_{kl}$, and
\[
  (v \wedge w)_{kl} = v_{k} w_{l} - v_{l} w_{k},
\]
which is the signed area of the projection of the parallelogram spanned by
$v$ and $w$ onto the plane of $e_{k}$ and $e_{l}$. Every bivector is a sum of
wedge products, but for $N \geq 4$ not every bivector is a single wedge
product: $v \wedge w$ has rank at most $2$, while $e_{12} + e_{34}$ has
rank $4$.

We use the inner product
$\langle X, Y\rangle = \frac12\operatorname{tr}(X^{\mathsf{T}}Y)
= \sum_{k<l} X_{kl}Y_{kl}$ on $\Lambda^{2}\mathbb{R}^{N}$, and
$|X| = \langle X, X\rangle^{1/2}$; the factor $\frac12$ accounts for each pair
$k < l$ appearing twice in $\operatorname{tr}(X^{\mathsf{T}}Y) = \sum_{i,j}
X_{ij}Y_{ij}$. Expanding the trace with
$\operatorname{tr}\bigl((pq^{\mathsf{T}})^{\mathsf{T}} rs^{\mathsf{T}}\bigr)
= (p \cdot r)(q \cdot s)$ gives
\[
  \langle a \wedge b, c \wedge d\rangle = (a \cdot c)(b \cdot d) - (a \cdot d)(b \cdot c).
\]
By this formula, $u_{k} \wedge u_{l}$, $k < l$, are orthonormal for every
orthonormal basis $u_{1}, \ldots, u_{N}$ of $\mathbb{R}^{N}$. Taking $c = a$
and $d = b$ in the formula gives the following.

\begin{lemma}[Lagrange's identity]
  \label{lem:lagrange}
  For $a, b \in \mathbb{R}^{N}$,
  $|a \wedge b|^{2} = |a|^{2}|b|^{2} - (a \cdot b)^{2}$.
\end{lemma}

We use Lemma~\ref{lem:lagrange} in the form $|a \wedge b| \leq |a|\,|b| \leq
\frac12(|a|^{2} + |b|^{2})$, where both inequalities are equalities if and only
if $a \cdot b = 0$ and $|a| = |b|$.

Let $\Pi \subseteq \mathbb{R}^{N}$ be a $2$-plane through the origin and
$\Pi^{\perp}$ its orthogonal complement. For an orthonormal basis $u_{1},
u_{2}$ of $\Pi$, the \emph{unit area element} of $\Pi$ is
$\omega_{\Pi} = u_{1} \wedge u_{2}$; for instance, $e_{kl}$ is that of the
plane spanned by $e_{k}$ and $e_{l}$. By bilinearity and antisymmetry,
changing the orthonormal basis multiplies $u_{1} \wedge u_{2}$ by the
determinant of the change of basis, which is $\pm 1$; statements about
$\omega_{\Pi}$ below hold for either sign. By Lemma~\ref{lem:lagrange}, $|\omega_{\Pi}| = 1$.

\subsection{The Hodge star on \texorpdfstring{$\Lambda^{2}\mathbb{R}^{4}$}{Lambda2R4}}

We use the Hodge star only for $N = 4$. For $k < l$ in $\{1, 2, 3, 4\}$ let
$\{k', l'\}$, $k' < l'$, be the complement of $\{k, l\}$, so that the plane
spanned by $e_{k'}, e_{l'}$ is the orthogonal complement of the plane spanned
by $e_{k}, e_{l}$. Let $\varepsilon_{kl} = \pm 1$ be the sign of the
permutation $(k, l, k', l')$ of $(1, 2, 3, 4)$. The \emph{Hodge star} is the
linear map $\star$ of $\Lambda^{2}\mathbb{R}^{4}$ with
\[
  \star e_{kl} = \varepsilon_{kl}\, e_{k'l'},
\]
that is, it sends the unit area element of each coordinate plane to that of its
orthogonal complement, up to sign:
\[
  \star e_{12} = e_{34}, \quad \star e_{34} = e_{12}, \quad
  \star e_{14} = e_{23}, \quad \star e_{23} = e_{14}, \quad
  \star e_{13} = -e_{24}, \quad \star e_{24} = -e_{13}.
\]
It permutes the orthonormal basis $e_{kl}$ up to signs, so it is an isometry.
Its use here is that, by Laplace expansion, it turns a $4 \times 4$
determinant into an inner product of bivectors.

\subsection{The Laplace expansion}

\begin{theorem}[{The Laplace expansion theorem, \cite[Section~0.8.9]{horn-johnson-2013}}]
  \label{thm:laplace}
  Let $A \in M_{n}(F)$, let $k \in \{1, \ldots, n\}$ be given, and let
  $\beta \subseteq \{1, \ldots, n\}$ be any given index set of cardinality $k$.
  Then
  \[
    \det A = \sum_{\alpha} (-1)^{p(\alpha,\beta)} \det A[\alpha, \beta]
      \det A[\alpha^{c}, \beta^{c}]
    = \sum_{\alpha} (-1)^{p(\alpha,\beta)} \det A[\beta, \alpha]
      \det A[\beta^{c}, \alpha^{c}]
  \]
  in which the sums are over all index sets $\alpha \subseteq \{1, \ldots, n\}$
  of cardinality $k$, and
  $p(\alpha, \beta) = \sum_{i \in \alpha} i + \sum_{j \in \beta} j$.
\end{theorem}

Here $M_{n}(F)$ denotes the $n \times n$ matrices over a field $F$,
$A[\alpha, \beta]$ is the submatrix of $A$ with rows indexed by $\alpha$ and
columns indexed by $\beta$, and $\alpha^{c}$, $\beta^{c}$ are the complements
of $\alpha$, $\beta$ \cite[Section~0.7.1]{horn-johnson-2013}. We use the first
expansion of Theorem~\ref{thm:laplace} for the $4 \times 4$ matrix $A$ with columns $v, v', w, w'$, index
sets of cardinality $2$, and $\beta = \{1, 2\}$, the columns of $v, v'$. For
$\alpha = \{k, l\}$, $k < l$, with complement $\{k', l'\}$, $k' < l'$, the
minor $\det A[\alpha, \beta]$ is $v_{k}v'_{l} - v_{l}v'_{k} =
(v \wedge v')_{kl}$, the complementary minor $\det A[\alpha^{c}, \beta^{c}]$
is $(w \wedge w')_{k'l'}$, and $(-1)^{p(\alpha, \beta)} = (-1)^{k + l + 1 + 2}$,
which is the sign $\varepsilon_{kl}$ of the permutation $(k, l, k', l')$.
Hence
\[
  \det(v, v', w, w')
  = \sum_{k<l} \varepsilon_{kl}\,(v \wedge v')_{kl}\,(w \wedge w')_{k'l'}.
\]
Since exchanging the two pairs is an even permutation,
$\varepsilon_{k'l'} = \varepsilon_{kl}$, so the coefficient of $e_{kl}$ in
$\star Y$ is $\varepsilon_{kl} Y_{k'l'}$, and this sum is an inner product:
\begin{equation}
  \label{eq:detstar}
  \det(v, v', w, w') = \langle v \wedge v', \star(w \wedge w')\rangle.
\end{equation}
For example, $\det(e_{1}, e_{2}, e_{3}, e_{4}) = \langle e_{12}, \star
e_{34}\rangle = \langle e_{12}, e_{12}\rangle = 1$. So, up to the areas
$|v \wedge v'|$ and $|w \wedge w'|$, the determinant compares the plane of
$v, v'$ with the orthogonal complement of that of $w, w'$.

\begin{remark}
  The description of $\star$ on coordinate planes holds for every plane: for
  a $2$-plane $\Pi \subseteq \mathbb{R}^{4}$ through the origin,
  \begin{equation}
    \label{eq:starplane}
    \star\,\omega_{\Pi} = \pm\,\omega_{\Pi^{\perp}}.
  \end{equation}
  Indeed, extend an orthonormal basis $u_{1}, u_{2}$ of $\Pi$ by one $u_{3},
  u_{4}$ of $\Pi^{\perp}$. By \eqref{eq:detstar}, $\langle u_{i} \wedge u_{j},
  \star(u_{1} \wedge u_{2})\rangle = \det(u_{i}, u_{j}, u_{1}, u_{2})$, which
  vanishes unless $\{i, j\} = \{3, 4\}$ and is $\pm 1$ then. The six
  $u_{i} \wedge u_{j}$, $i < j$, are orthonormal, hence a basis of
  $\Lambda^{2}\mathbb{R}^{4}$, and these inner products are the coordinates
  of $\star\omega_{\Pi}$ in it; so $\star\omega_{\Pi} = \pm\, u_{3} \wedge
  u_{4} = \pm\,\omega_{\Pi^{\perp}}$.
\end{remark}
